\section{Training a stronger Moving initial forecast}
We compare two objectives for the same CV-anchored velocity spline
with a bounded amplitude. The control maximizes conservative gain
against OLS4, debiting forecast disagreement at missing future steps.
The second objective maximizes the worse of the CV- and OLS4-relative
gains and penalizes positive excess over the original all, easy and
signed CV risks, with a soft incentive for reference support.
Each minibatch belongs to one fitting recording. The utility loss
weights queries equally while the risk surrogate uses unweighted
rows. This is a fixed training penalty with separate calibration
sources~\cite{cotter}, not a constrained optimization or a coverage
guarantee. Both models read only observed trajectories and their
causal neighbors.

We train \MovingInitialFits{} models for \MovingInitialUpdates{}
updates, with three seeds and disjoint fitting, calibration and
held sources per rotation. The \MovingInitialViews{} views represent
\MovingInitialRecords{} recordings from \MovingInitialSources{}
already used development sources, not independent test data.
A strict choice must pass the unchanged 2\% rule in every
calibration recording and give positive source-macro conservative
gain against both references. A diagnostic choice ranks by joint
passes, then by the smaller conservative gain. EqMotion and PCT are
held-only controls because their existing calibration forecasts
are not proven to exclude the outer held sources. We freeze all
predictions before held scoring. No strict choice qualifies
(\MovingInitialStrict{}/\MovingInitialCells{}); CV fallback
does not turn undefined support into a pass.

\begin{center}
\small
\setlength{\tabcolsep}{3.5pt}
\begin{tabular}{lrrr}
\toprule
Initial plan & Risk & Dual & Matched dual \\
\midrule
Half-peak anchor & \MovingInitialAnchorRisk & \MovingInitialAnchorDual & \MovingInitialAnchorMatchedDual \\
OLS4 objective & \MovingInitialOLSRisk & \MovingInitialOLSDual & \MovingInitialOLSMatchedDual \\
Dual-risk objective & \MovingInitialRiskRisk & \MovingInitialRiskDual & \MovingInitialRiskMatchedDual \\
Diagnostic choice & \MovingInitialScreenRisk & \MovingInitialScreenDual & \MovingInitialScreenMatchedDual \\
\bottomrule
\end{tabular}
\end{center}
Columns count risk passes, joint risk and utility passes, and
joint passes with the old budget, respectively.

The OLS4 objective improves on OLS4 where labels are available
(\MovingInitialOLSKnown{} pp), but loses with missing-step debits
(\MovingInitialOLSLower{} pp). The dual-risk objective also loses
after debits (\MovingInitialRiskLower{} pp); its smaller number
of edited rows does not imply supported risk. Unmatched diagnostic
paths narrow the relaxed single-query exclusion count from
\MovingInitialAnchorUpper{} to \MovingInitialRawUpper{}/\MovingInitialViews{}
views, but use their own intervention budget and do not add a dual pass.

\paragraph{Fixed recipients and peaks.}
An attribution audit locks all \MovingInitialEdits{} old edits,
including their query times, recipients and individual peaks. We
take the new direction only when it reaches that peak without
amplification; otherwise the old path stays. Against the old anchor,
the OLS4 objective gives \MovingInitialOLSMatched{} pp, the dual-risk
objective gives \MovingInitialRiskMatched{} pp, and the diagnostic
choice gives \MovingInitialScreenMatched{} pp in conservative gain.
Brackets give 95\% intervals from \MovingInitialBootstrap{} paired
source draws after averaging seeds. This does not establish a better
direction at fixed budget or justify deployment. The original
risk rule and the independent confirmation boundary are unchanged.
